\documentclass[11pt,a4paper]{article}

% ----------------------- Packages -----------------------
\usepackage[utf8]{inputenc}
\usepackage[T1]{fontenc}
\usepackage{geometry}
\geometry{margin=2.5cm}
\usepackage[pdftex,hidelinks]{hyperref}
\usepackage{enumitem}
\usepackage{titlesec}

% ----------------------- Heading styling: plain, no colour -----------------------
\titleformat{\section}{\Large\bfseries}{\thesection.}{0.7em}{}
\titleformat{\subsection}{\normalsize\bfseries}{\thesubsection.}{0.7em}{}
\titlespacing*{\section}{0pt}{1.2em}{0.5em}
\titlespacing*{\subsection}{0pt}{0.9em}{0.3em}

% ----------------------- Title / author -----------------------
\title{\vspace{-1.5cm}\Huge\textbf{FundiLink}\\[0.5em]
\Large Concept Paper\\[0.3em]
\normalsize Connecting Communities to Trusted Local Artisans}
\author{Prepared for Prospective Partners and Funders}
\date{\today}

\begin{document}

\maketitle

% ======================= 1. The Bottom Line =======================
\section{The Bottom Line}

Skilled manual work -- plumbing, electrical repairs, carpentry, painting -- is essential to
every home and small business. Yet in our market, getting this kind of help is a constant
gamble: you do not know who is qualified, what the job should cost, or whether the work will
be done properly.

\textbf{FundiLink solves this by bringing order and trust to an informal market.} At its core
it is a location-based service: using your location, it shows the qualified artisans working
near you on a map, and you choose the one you want. It sets a clear agreed price up front and
builds each artisan's reputation through genuine, public reviews.

Think of it like the way a trusted ride-hailing service shows you available drivers near you on
a map, and you pick the one you want -- but for skilled home services. Instead of waiting days
and taking a chance on whoever you can find, you see the vetted artisans around you and choose
in minutes.

% ======================= 2. The Problem =======================
\section{The Problem}

Finding a good artisan today is difficult. The reasons are straightforward.

\subsection{You Cannot Verify Who Is Qualified}

There is no reliable way to tell a skilled, trustworthy artisan from an unreliable one before
letting them into your home. No directory or registry exists, and qualifications are rarely
verifiable. A customer must simply take a chance.

\subsection{Prices Are Unclear and Unfair}

Quotes are almost always verbal. For the same job, different artisans may quote wildly
different prices, and the final bill often ends up higher than the original promise. Without
anything in writing, the customer has little ground to dispute it.

\subsection{Reputation Does Not Exist}

A skilled artisan's past performance is invisible. A person who does excellent work cannot
build a reputation or attract new customers on merit, while a person who does poor work faces
no consequence. There is no incentive for quality, and no way for customers to learn from
others' experiences.

\subsection{Finding Help Quickly Is Hard}

When a repair is urgent -- a burst pipe, a failed electrical circuit -- there is no fast,
reliable way to see which qualified professionals are nearby and available. People often wait
days because they simply cannot find anyone close to help.

Taken together, these problems mean a service everyone needs is delivered slowly, unfairly and
unreliably.

% ======================= 3. Why Now =======================
\section{Why Now}

This is the right time to solve this problem, because the ingredients for change are already
in place:

\begin{itemize}[leftmargin=1.5em, label={--}]
    \item \textbf{Growing cities.} More homes are being built, bought and maintained, and
          demand for reliable repair and improvement services is rising quickly.
    \item \textbf{More people on phones.} A large and growing share of households now use
          smartphones, so a phone-based solution can reach them easily.
    \item \textbf{A proven approach.} The idea of matching customers with nearby providers and
          publishing honest reviews has transformed other sectors. It has simply not yet been
          applied to skilled home services here.
\end{itemize}

FundiLink brings these already-present habits together. We are not asking people to change
their behaviour; we are taking how they already hire artisans and making it faster, fairer and
more reliable.

% ======================= 4. How We Solve It =======================
\section{How We Solve It}

The heart of the solution is simple: turn an informal, uncertain handshake into a clear,
ordered process. Each job on FundiLink follows the same sequence.

\begin{enumerate}[leftmargin=1.5em]
    \item \textbf{Find on the map.} Using the customer's location, the map shows the qualified
          artisans working nearby, with their ratings and the work they do.
    \item \textbf{Choose a fundi.} The customer picks the fundi they want from the map and
          sends them a request.
    \item \textbf{Agreed price.} Customer and fundi agree on a clear price before any work
          begins, so there is no surprise on the final bill.
    \item \textbf{Watching the job.} The customer can see when the fundi is on the way, has
          arrived, and is working, so there is no wondering or waiting in the dark.
    \item \textbf{Completion and review.} Once the customer confirms the work is done, they
          leave a rating and a review that help the next customer decide.
    \item \textbf{Fairness, always.} If anything goes wrong, a neutral team can step in, review
          what happened, and settle the matter fairly.
\end{enumerate}

\subsection{How Trust Is Built and Kept}

The greatest long-term value of FundiLink is the reputation it creates. Every completed job
adds to a public record of how each artisan performs. Over time, reliable, skilled artisans
are rewarded with more work and higher standing, and those who do poor work naturally lose
business. Customers gain a growing, trustworthy list of professionals they can rely on.

This is what turns a risky, uncertain market into one where quality is rewarded -- and it is
the strongest force for good service available.

\subsection{Designed Around the Customer}

FundiLink is built for how people really live and work here: signing up with a phone number,
communicating in a local language, and using a simple experience. The whole journey is made
easy, because the point is to take the stress out of getting help, not to add complexity.

% ======================= 5. The MVP =======================
\section{Our First Step: A Clear, Focused MVP}

The full vision is ambitious, but our first version does not need to be. A Minimum Viable
Product (MVP) delivers the one thing that matters most -- getting a customer a trusted
artisan, nearby, quickly -- and nothing more. Starting small means we can launch sooner, learn
from real use, and improve based on what customers and artisans actually need.

The MVP focuses on a single city and a narrow set of essential features, each chosen to prove
the core idea end to end.

\subsection{The Core: Finding Fundis on the Map}

The heart of the MVP is geolocation. Using the customer's location, the map instantly shows the
qualified fundis working nearby -- much like a trusted ride-hailing service shows the drivers
around you. The customer picks the fundi they want from the map. This turns a day-long search
into a quick choice of someone near enough to help.

\subsection{Essential Features in the First Version}

\begin{itemize}[leftmargin=1.5em, label={--}]
    \item \textbf{Simple sign-up.} Customers and fundis join using a phone number, with no
          complicated forms or paperwork.
    \item \textbf{See fundis on the map.} Using geolocation, the map shows the qualified fundis
          working nearby, with their ratings and the work they do.
    \item \textbf{Choose and request.} The customer picks a fundi from the map, and with a few
          taps sends them a request describing the job.
    \item \textbf{Agreeing on a price.} Customer and fundi agree on a clear price before work
          starts, removing surprise on the final bill.
    \item \textbf{Rating and review.} After each job, the customer rates the work, building a
          store of honest, public feedback that guides future customers.
\end{itemize}

\subsection{What We Deliberately Leave Out -- for Now}

The MVP keeps things simple by omitting features that can come later, including extra service
categories, advanced administrative tools, and expansion beyond the first city. By focusing
only on the essentials, we can get the trusted, nearby-matching experience into people's hands
quickly and make it excellent before we grow.

% ======================= 6. What We Need =======================
\section{What We Need to Make It Happen}

FundiLink is already working as a proven prototype. The whole journey -- seeing fundis on the
map, choosing one, agreeing on a price, following the job, and leaving a review -- has been
built and demonstrated from start to finish. Support will take it from this working foundation
to a trusted service people rely on daily. We will focus on a few things.

\subsection{Perfecting the Core Experience}

The most important step is to make the core journey excellent: a fast, reliable map showing
nearby fundis, clear agreed pricing, and an easy way to follow the job. We will invest in
making this experience dependable and polished for every customer and fundi.

\subsection{Launching in Our First City}

We will launch in one city and do it properly: recruiting and checking a strong base of
skilled artisans, building an initial base of customers, and running the activities needed to
reach the point where the marketplace keeps growing on its own.

\subsection{Building Our Team}

We are a small team today. Support will let us grow -- bringing on the people needed to build,
support and improve the service, and to keep it reliable for every user.

\subsection{Strengthening and Growing the Service}

We will keep improving: deeper local-language support, stronger protections for both customers
and artisans, and better tools for running the marketplace fairly and well as it grows.

% ======================= 7. Conclusion =======================
\section{Conclusion}

There is a large, urgent and unmet need for a trustworthy way to find, hire and review skilled
artisans. The habits that make such a service possible -- growing phone use and comfort with
booking services -- are already present. What is missing is a service that brings them
together for skilled manual work.

FundiLink is that service. It has been designed, built and proven, and it is ready to be taken
fully to market with the right support. We would welcome the opportunity to discuss the
project with you, to show the working service, and to explore a partnership together.

\end{document}
